This thesis examines how the super-sample effect and the box replication effect influence the accuracy and interpretation of higher-order statistics from weak gravitational lensing data. These effects are critical in understanding the limitations and potential biases in simulations used for cosmological studies. The study begins with a review of modern cosmology. It highlights how weak lensing measures small distortions in the shapes of distant galaxies caused by intervening matter. This can help us understand dark matter, dark energy, and the overall structure of the Universe. It then focuses on beyond-two-point statistics, a method reveals non-Gaussian features in the matter distribution and thereby enhance constraints on cosmological parameters.

To achieve this, we use a simulation-based framework with two $N$-body simulation strategies: the BIGBOX simulation set from the HalfDome project, which has $6144^3$ particles in a $3750\, \mathrm{Mpc}/h$ box, and the TILED simulation set, composed of multiple $1024^3$-particle boxes, each $625\, \mathrm{Mpc}/h$ in size. Both simulations share the same cosmological parameters but differ in volume coverage and box replication schemes, allowing us to isolate super-sample modes and other effects. We project three-dimensional matter distributions into full-sky convergence maps at various redshifts under the Born approximation, then incorporate shape noise and Gaussian smoothing to approximate realistic observing conditions. From these maps, we employ a Fibonacci grid to extract numerous $10^\circ \times 10^\circ$ patches and evaluate a comprehensive set of lensing observables—namely, the angular power spectrum, bispectrum, probability distribution function (PDF), peak and minima counts, and Minkowski functionals. Comparisons of the mean values, variances, and covariances of these observables between BIGBOX and TILED are further contextualized by theoretical models such as \texttt{Halofit} and \texttt{hmpdf}.

\at{Our findings show that the mean values of all investigated statistics from BIGBOX and TILED agree at the $1\%$ level, confirming consistency under the same cosmological framework. However, variance and covariance estimates differ by up to $10\%$-$30\%$ at higher multipoles and redshifts, primarily due to super-sample modes present in the larger-volume BIGBOX simulations but absent in TILED. Although shape noise and Gaussian smoothing affect certain observables, most notably the angular power spectrum and the counts of peaks and minima, the super-sample effect remains the dominant contributor to these discrepancies. Furthermore, we detect $2\%$-$5\%$ systematic biases introduced by box replication artifacts, which can obscure super-sample signals by shifting the means of selected observables and altering the distribution of prominent peaks and deep minima. These results highlight several key considerations for future research and the design of upcoming weak lensing surveys.} 

\at{First, large-scale modes must be incorporated into simulations to ensure realistic uncertainty estimates. We adopt the approximate dimensional analysis}
$
R_{\sigma^2} \propto \left(L^2 / \Omega_W\right)^{\alpha},
$
\at{where $L$ is the simulation box size, $\Omega_W$ is the survey area, and $\alpha$ is an undetermined exponent. This implies that super-sample variance becomes increasingly significant for larger surveys employing relatively small simulation boxes. For instance, matching the variance levels of TILED in an $18{,}000\ \mathrm{\deg}^2$ survey (e.g., LSST) would require a simulation box of approximately $8\ \mathrm{Gpc}/h$.}

\at{Secondly, shape noise and Gaussian smoothing must be carefully managed to maintain the fidelity of higher-order statistics. Although these effects can shift certain observables, the super-sample effect remains robust under various observational conditions.}

\at{Finally, mitigating box replication artifacts is essential for high-precision analyses. These artifacts introduce $2\%$-$5\%$ biases in mean observables and can modify covariance structures by up to $10\%$. Nonetheless, improved simulation and data analysis strategies can address these systematic issues.}

\at{This work highlight the importance of quantifying super-sample effects, particularly in higher-order statistics, an area where no prior work existed. Additionaly, it provides one of the first comprehensive investigations into the impact of box replication effect on summary statistics and their covariance. As forthcoming surveys, including the Rubin Observatory's LSST, Euclid, and the Roman Space Telescope, aim for increasingly stringent constraints, properly accounting for super-sample effects and replication artifacts will be indispensable. Future efforts will focus on reducing replication biases, exploring additional cosmological models and baryonic physics, and integrating these advances into survey planning and analysis pipelines. By refining both simulation methodologies and interpretive frameworks, this thesis helps establish weak lensing as an even more robust and informative probe of the Universe's fundamental properties.}
